\section{Continuation complexity of a finite-data test}
\label{sec:v19-streaming}
The reset budget does not by itself bound the auditor's retained memory.
We now impose both constraints. At scheduled epoch $t$ a streaming auditor
reads one letter $x_t$ from a finite alphabet $X_t$, updates a register,
and emits no information to any auxiliary store. It makes its decision on
reading the last letter. Transition rows may depend on the epoch; the clock
is free. All other information crossing a cut belongs to the register.
Fresh randomization is allowed. An independent persistent selector, when
allowed, selects a complete auditor before any letters are read.

For a Boolean decision $\psi:X_1\times\cdots\times X_n\to\{-1,1\}$,
let $u\sim_t v$ mean that
\[
 \psi(us)=\psi(vs)\quad\hbox{for every }s\in X_{t+1}\times\cdots\times X_n.
\]
Write $r_t(\psi)$ for the number of equivalence classes of prefixes of
length $t$, for $0\le t<n$. These are the deterministic continuation
responses. The terminal decision is emitted on the last transition, so no
terminal register is needed. In particular, $r_0=1$. The equivalence is a
finite-horizon version of the classical residual-state construction
\cite{Nerode1958}. The following quantitative statement will also measure
the loss when a statistical auditor has fewer states than this construction.

\begin{theorem}[Exact continuation width and a quantitative cut obstruction]
\label{thm:v19-residual}
A stochastic streaming auditor computes $\psi$ exactly on every word with
width profile $(K_t)_{0\le t<n}$ if and only if
$K_t\ge r_t(\psi)$ at every cut. The profiles are attained simultaneously
by a deterministic auditor. Independent persistent randomization does not
reduce these widths for exact computation.

More generally, at a cut with $r=r_t(\psi)\ge2$, every stochastic auditor
of width $K$ satisfies
\begin{equation}\label{eq:v19-cuterror}
 \sum_{w\in X_1\times\cdots\times X_n}
       e_\psi(w)\ \ge\ \frac{(r-K)_+}{r-1},
\end{equation}
where $e_\psi(w)$ is its probability of an incorrect Boolean decision on
input $w$. The inequality also holds after averaging over an independent
selector whose conditional auditors each have width at most $K$.
\end{theorem}
\begin{proof}
From a prefix $u$ let $e_u(m)$ be the distribution of the retained state.
From state $m$ and suffix $s$ let $a_m(s)$ be the probability of terminal
answer $+1$. Exact computation says
\[
 \sum_m e_u(m)a_m(s)=\frac{1+\psi(us)}2.
\]
The right side is zero or one. Every state of positive $e_u$-mass must
therefore have precisely that zero--one continuation response for all $s$.
Distinct prefix classes have disjoint positive supports. Hence $K_t\ge r_t$.
Conversely, retain the prefix class. Appending the same letter to two
equivalent prefixes gives equivalent new prefixes; thus every transition
is well defined. At the last transition the class and the last letter
determine $\psi$. This realizes every width $r_t$ simultaneously.
For a persistent independent selector, a zero average error on each of the
finitely many words implies zero conditional error on all words outside a
single null set of selector values. The preceding bound applies there.

For the quantitative claim choose representative prefixes $u_1,\ldots,u_r$,
one per class, and put $e_i(m)=e_{u_i}(m)$. Define
\[
 o_{ij}=\sum_{m=1}^K\min\{e_i(m),e_j(m)\}.
\]
For nonnegative numbers $a_1,\ldots,a_r$, ordering them shows that
$\sum_{i<j}\min(a_i,a_j)\ge\sum_i a_i-\max_i a_i$.
Apply this at each state. Since every $e_i$ has mass one and each
$\max_i e_i(m)\le1$, it follows that
\begin{equation}\label{eq:v19-overlap}
 \sum_{i<j}o_{ij}\ge r-K.
\end{equation}
Choose a suffix $s_{ij}$ distinguishing $u_i,u_j$. Conditional on a common
state and this suffix, the two desired answers are opposite. Their error
probabilities, weighted by $e_i(m),e_j(m)$, have sum at least
$\min\{e_i(m),e_j(m)\}$. Thus
\[
 e_\psi(u_i s_{ij})+e_\psi(u_j s_{ij})\ge o_{ij}.
\]
On summing, each complete word on the left occurs at most $r-1$ times:
its prefix identifies a single representative, which participates in
$r-1$ pairs. Nonnegativity of all other word errors and
\eqref{eq:v19-overlap} prove \eqref{eq:v19-cuterror}. When $r\le K$ the
claim is just nonnegativity. Finally apply the bound conditional on a
selector and integrate. Independence ensures that conditioning has not
changed the input experiment.
\end{proof}

The exact part is also the extremal zero--one case of
Theorem~\ref{thm:v18-continuation}. Its content is a simultaneous width
profile, rather than separate unrelated factorizations at different cuts.
The error bound does not assert that its constant is optimal for every
Boolean function. Its use below is to convert a continuation obstruction
into a strict loss of the attainable statistical value.

\subsection{The width required to attain the binary reset value}
Fix $0<h<1/2$ and the binary candidate class $B_h$ of
Theorem~\ref{thm:v18-auditlower}. A terminal decision $+1$ chooses the
singleton event appropriate to the $+$ candidate, and $-1$ the complementary
event. A zero-score event can equivalently be replaced by an independent
fair choice of these decisions. The expected payoffs are therefore
$h\E_+d$ and $-h\E_-d$, where $d\in[-1,1]$ is the conditional mean
decision and $P_\pm=\operatorname{Ber}(1/2\pm h)^{\otimes n}$.
Let $v_{n,K}$ be the supremum of the smaller of these two payoffs over
streaming auditors with the declared width profile. The final validation
interface is unchanged; the widths here count training and selection only.

\begin{theorem}[Exact memory threshold for optimal binary auditing]
\label{thm:v19-binarywidth}
Let $n=2m+1$ be odd. The unconstrained reset value is
$v_n=h\|P_+-P_-\|_{\rm TV}$. It is attained at the profile
\begin{equation}\label{eq:v19-majorityprofile}
 r_t=\min\{t+1,n-t+2\},\qquad 0\le t<n,
\end{equation}
and $v_{n,K}=v_n$ if and only if $K_t\ge r_t$ for all these cuts.
For every cut with $r_t\ge2$ one has the explicit deficit
\begin{equation}\label{eq:v19-memorydeficit}
 v_n-v_{n,K}\ \ge\
 2h^2(1/4-h^2)^m\frac{(r_t-K_t)_+}{r_t-1}.
\end{equation}
These conclusions remain valid when an independent persistent selector
chooses auditors of the same conditional widths. For $n\ge3$ the least
peak training width attaining $v_n$ is $m+2$; for $n=1$ it is one.
\end{theorem}
\begin{proof}
The sign of $P_+(w)-P_-(w)$ is the majority decision $\psi(w)$, with no
zero difference since $n$ is odd. Complementing every bit interchanges the
two laws and reverses the decision. Consequently majority has equal
payoffs and attains $h\|P_+-P_-\|_{\rm TV}$, as also proved in
Theorem~\ref{thm:v18-auditlower}.

A prefix with $k$ ones requires at least $m+1-k$ ones in its remaining
$n-t$ places. Thresholds at most zero give the constant positive response;
thresholds greater than $n-t$ give the constant negative response; the
remaining thresholds give different nonconstant responses. Counting them
for $0\le k\le t$ gives \eqref{eq:v19-majorityprofile}. The exact
realization in Theorem~\ref{thm:v19-residual} implements majority at all
these widths.

For an arbitrary randomized decision, its smaller payoff is bounded by
the equal-prior average. The deficit of this average from the majority
value is exactly
\begin{equation}\label{eq:v19-weightedregret}
 h\sum_w |P_+(w)-P_-(w)|e_\psi(w).
\end{equation}
Put $a=1/2+h$ and $b=1/2-h$. A word with $k\ge m+1$ ones has absolute
likelihood difference
$(ab)^{n/2}\{(a/b)^{k-n/2}-(a/b)^{-(k-n/2)}\}$, which increases with
$k$. Symmetry treats the other half of the words. The smallest difference
is therefore attained at $k=m,m+1$ and equals
\[
 a^{m+1}b^m-b^{m+1}a^m=2h(1/4-h^2)^m.
\]
Substitution of \eqref{eq:v19-cuterror} into
\eqref{eq:v19-weightedregret} gives \eqref{eq:v19-memorydeficit}.
Any deficient cut gives a strictly positive loss; all sufficiently wide
profiles implement majority. Maximizing \eqref{eq:v19-majorityprofile}
proves the peak-width assertion. The same average-payoff argument works
conditional on an independent selector and hence after averaging.
\end{proof}

The result identifies the exact memory needed to attain the optimal value,
not the memory needed for a prescribed approximate value uniformly in $h$.
The quantitative deficit can be small for large $n$. In particular it is
not a sharp joint asymptotic tradeoff in sample size, state size and error.
It complements, rather than replaces, the uniform sampling exponent in
Theorem~\ref{thm:v18-auditlower}. No minimax interchange on the generally
nonconvex set of fixed-width auditors has been used.
